% Complete Section 7.5, original PDF 98-100 / printed 88-90.
% block: finite-element-space-properties-heading page:98 type:section
\subsection{{Properties of the finite element spaces}}\label{sec:finite-element-space-properties}

% block: finite-element-approximation-introduction page:98 type:ordinary-paragraph
{The approximation properties of the finite element spaces $\mathcal P_r\Lambda^k(\mathcal T_h)$ and $\mathcal P_r^-\Lambda^k(\mathcal T_h)$ (asymptotically, as the mesh size $h$ tends to $0$) can be established as is usual for finite elements, using the fact that the canonical projection preserves polynomials locally, and applying the Bramble–Hilbert lemma and scaling. This is carried out in detail in [\BookNumber{11}, Theorem 5.3], where the following result is proved.}

% block: canonical-projection-approximation-theorem page:98 type:theorem
\Needspace{8\baselineskip}
\begin{Theorem}\label{thm:canonical-projection-approximation}
{Denote by $\Pi_h$ the canonical projection of $\Lambda^k(\Omega)$ onto either $\mathcal P_r\Lambda^k(\mathcal T_h)$ or $\mathcal P_{r+1}^-\Lambda^k(\mathcal T_h)$. Let $1\leq p\leq\infty$ and $(n-k)/p<s\leq r+1$. Then $\Pi_h$ extends boundedly to $W_p^s\Lambda^k(\Omega)$, and there exists a constant $C$ independent of $h$, such that}
\[
\|\omega-\Pi_h\omega\|_{L_p\Lambda^k(\Omega)}\leq Ch^s|\omega|_{W_p^s\Lambda^k(\Omega)},\qquad\omega\in W_p^s\Lambda^k(\Omega).
\]
\end{Theorem}

% block: canonical-projection-trace-restriction page:98 type:ordinary-paragraph
{Note that the restriction $s>(n-k)/p$ comes from the Sobolev embedding theorem: it is what is required for functions in $W_p^s$ to having traces on $k$-dimensional faces.}

% block: complete-and-trimmed-form-errors pages:98,99 type:ordinary-paragraph
{For a smooth $k$-form $\omega$ and the complete polynomial space $\mathcal P_r\Lambda^k(\mathcal T_h)$, we may take $p=2$, $s=r+1$ to obtain an $L^2$ estimate of order $r+1$:}
\[
\|\omega-\Pi_{r,h}^k\omega\|_{L^2\Lambda^k}\leq Ch^{r+1}|\omega|_{H^{r+1}\Lambda^k}.
\]
{For the trimmed polynomial space $\mathcal P_r^-\Lambda^k(\mathcal T_h)$, which contains the full polynomial space $\mathcal P_{r-1}\Lambda^k(\mathcal T_h)$ but not $\mathcal P_r\Lambda^k(\mathcal T_h)$, the order is reduced to $r$:}
\[
\|\omega-\Pi_{r,h}^{-,k}\omega\|_{L^2\Lambda^k}\leq Ch^r|\omega|_{H^r\Lambda^k}.
\]

% block: complete-form-derivative-error page:99 type:ordinary-paragraph
{For the complete polynomial space, we have the cochain relation}
\[
d\Pi_{r,h}^k\omega=\Pi_{r-1,h}^{k+1}d\omega,
\]
{and so}
\[
\|d(\omega-\Pi_{r,h}^k\omega)\|_{L^2\Lambda^{k+1}}=\|d\omega-\Pi_{r-1,h}^{k+1}d\omega\|_{L^2\Lambda^{k+1}}\leq Ch^r|d\omega|_{H^r\Lambda^{k+1}},
\]
{a loss of one order of convergence for $d\omega$ as compared to $\omega$. For the trimmed space the cochain relation is}
\[
d\Pi_{r,h}^{-,k}\omega=\Pi_{r,h}^{-,k+1}d\omega,
\]
{so there is no loss of order:}
\[
\|d(\omega-\Pi_{r,h}^{-,k}\omega)\|_{L^2\Lambda^{k+1}}\leq Ch^r|d\omega|_{H^r\Lambda^{k+1}}.
\]

% block: four-finite-element-pairs page:99 type:ordinary-paragraph
{We are now ready to choose the pair of Galerkin subspaces $V_h^{k-1}$, $V_h^k$. Two obvious choices are}
\[
\mathcal P_r\Lambda^{k-1}(\mathcal T_h),\mathcal P_{r-1}\Lambda^k(\mathcal T_h)\qquad\text{and}\qquad\mathcal P_r^-\Lambda^{k-1}(\mathcal T_h),\mathcal P_r^-\Lambda^k(\mathcal T_h),
\]
{using either complete polynomial spaces for both subspaces or trimmed spaces for both. There are two other choices as well, which use one complete space and one trimmed. We may reduce $V_h^{k-1}$ in the complete polynomial pair to obtain}
\[
\mathcal P_r^-\Lambda^{k-1}(\mathcal T_h),\mathcal P_{r-1}\Lambda^k(\mathcal T_h).
\]
{This does not change the range $dV_h^{k-1}$ and again the canonical projections commute: $d\Pi_{r,h}^{-,k}=\Pi_{r-1,h}^kd$. Similarly, we may expand $V_h^{k-1}$ in the trimmed polynomial pair to get}
\[
\mathcal P_r\Lambda^{k-1}(\mathcal T_h),\mathcal P_r^-\Lambda^k(\mathcal T_h),
\]
{which, again, does not affect the range of $d$. Altogether, we obtain four possible choices for each degree $r$:}
\begin{equation}\label{eq:four-galerkin-form-pairs}
V_h^{k-1}=\left\{\begin{gathered}\mathcal P_r\Lambda^{k-1}(\mathcal T_h)\\\text{or}\\\mathcal P_r^-\Lambda^{k-1}(\mathcal T_h)\end{gathered}\right\},\qquad
V_h^k=\left\{\begin{gathered}\mathcal P_r^-\Lambda^k(\mathcal T_h)\\\text{or}\\\mathcal P_{r-1}\Lambda^k(\mathcal T_h)\ (\text{if }r>1)\end{gathered}\right\}.
\end{equation}

% block: galerkin-form-pair-error-rates page:99 type:ordinary-paragraph
{The $L^2$ rates of approximation they achieve for the variables $u\in\Lambda^k$ and $\sigma=\delta u\in\Lambda^{k-1}$ are shown in Table~\ref{tab:galerkin-form-approximation-rates}. That these estimates hold not only for the projection of $\sigma$ and $u$ but also for the Galerkin solution (when the true solution is smooth) can be deduced from the improved error estimates of Theorem~\ref{thm:mixed-improved-errors}.}

% block: third-galerkin-form-space page:99 type:ordinary-paragraph
{Of course, to obtain these results we need to fulfill the hypotheses of Theorem~\ref{thm:mixed-improved-errors}. For this we have to specify a third space $V_h^{k+1}\subset H\Lambda^{k+1}$. (This space is not used in the Galerkin method but is required for theory.) For this we take}
\[
V_h^{k+1}=\begin{cases}\mathcal P_r^-\Lambda^{k+1}(\mathcal T_h)&\text{if }V_h^k=\mathcal P_r^-\Lambda^k(\mathcal T_h),\\\mathcal P_{r-1}^-\Lambda^{k+1}(\mathcal T_h)&\text{if }V_h^k=\mathcal P_{r-1}\Lambda^k(\mathcal T_h).
\end{cases}
\]

% block: galerkin-form-approximation-rates-table page:100 type:table
\begin{table}[htbp]\centering
\caption{{Rates of approximation for stable choices of Galerkin subspaces.}}\label{tab:galerkin-form-approximation-rates}
\begin{tabular}{cc|cccc}\hline
$V_h^{k-1}$&$V_h^k$&$\sigma$&$d\sigma$&$u$&$du$\\\hline
$\mathcal P_r\Lambda^{k-1}$&$\mathcal P_r^-\Lambda^k$&$r+1$&$r$&$r$&$r$\\
$\mathcal P_r^-\Lambda^{k-1}$&$\mathcal P_r^-\Lambda^k$&$r$&$r$&$r$&$r$\\
$\mathcal P_r\Lambda^{k-1}$&$\mathcal P_{r-1}\Lambda^k$&$r+1$&$r$&$r$&$r-1$\\
$\mathcal P_r^-\Lambda^{k-1}$&$\mathcal P_{r-1}\Lambda^k$&$r$&$r$&$r$&$r-1$\\\hline
\end{tabular}
\end{table}

% block: bounded-form-cochain-projections page:100 type:ordinary-paragraph
{This defines the subcomplex \eqref{eq:galerkin-subcomplex} of the corresponding segment of the de Rham complex, required by Theorem~\ref{thm:mixed-improved-errors}. The remaining hypothesis is that there exist $W$-bounded cochain projections, i.e., bounded projection operators $\pi_h^j:L^2\Lambda^j\to V_h^j$ for $j=k-1$, $k$, $k+1$ which commute with $d$ on its domain. We have already remarked that the canonical projections, defined through the DOFs, form a cochain projection, that is, they commute with $d$ on smooth enough forms. However, the canonical projection is not in general bounded on $L^2\Lambda^k$, nor even on $H\Lambda^k$. For example, $\Pi_{1,h}^0$ is the usual interpolant into the space of piecewise linear functions. It is not bounded on $H\Lambda^0=H^1$ in more than one dimension, because it depends on the point values at the vertices, and it is not $L^2$ bounded even in one dimension. However, there are several ways a bounded cochain projection can be constructed. One approach, which is presented in detail in [\BookNumber{11}, Section 5.4], is to apply a mollifier before applying the canonical projection (following \BookCite{28}), and then to correct the result to obtain a projection with the desired properties (following \BookCite{66}). This construction provides the desired $W$-bounded cochain projections, and hence the error estimates of Theorems~\ref{thm:mixed-basic-error} and~\ref{thm:mixed-improved-errors} hold.}

% block: cochain-projection-further-properties page:100 type:ordinary-paragraph
{We remark that there are various other properties one might desire from the cochain projections which are not provided by the construction just recalled. For example, one might want the construction to be local or one might ask that the bounds be uniform in the degree $r$. For this reason, the construction of bounded cochain projections remains an active research area \BookCite{30}, \BookCite{43}.}
